\section{Proofs for Section~\ref{sec:known-parameter}}
\label{app:known-parameter}


\subsection{Fixed-Scale Inexact UTR}
\label{app:fixed-scale-proofs}

Let $K\in\mathbb{N}_0\cup\{\infty\}$ denote the first return index of
the procedure
\eqref{eq:fixed-scale-stopping}--\eqref{eq:fixed-scale-subproblem},
with $K=\infty$ if the procedure never returns.
For $k<K$, set
\[
    r_k:=\|d_k\|,
\]
and let $\lambda_k$ be the trust-region multiplier associated with
\eqref{eq:fixed-scale-subproblem}.
Thus Lemma~\ref{lem:tr-global-optimality} is applied with
\[
    B
    =
    H_k+\xi_kI
    +2\sqrt{\frac M6}\sqrt{h_k}\,I,
    \qquad
    R
    =
    \frac{\sqrt{h_k}}{4\sqrt{M/6}}.
\]


\begin{lemma}[Decrease and clipped-gradient recurrence]
\label{lem:fixed-one-step}
Under the assumptions of Theorem~\ref{thm:fixed-scale-utr},
every nonterminal step $k<K$ satisfies
\begin{equation}
    F(x_k)-F(x_{k+1})
    \ge
    \frac34
    \sqrt{\frac M6}\sqrt{h_k}\,r_k^2
    +
    \frac12\lambda_kr_k^2
    -
    \frac{\epsilon r_k}{8192}.
    \label{eq:fixed-basic-decrease}
\end{equation}
Consequently,
\begin{subequations}
\label{eq:fixed-step-type-decrease}
\begin{align}
    \lambda_k=0:\qquad
    F(x_k)-F(x_{k+1})
    &\ge
    -
    \frac{
        \epsilon^{3/2}
    }{
        201\,326\,592\sqrt{M/6}
    },
    \label{eq:fixed-interior-increase}
    \\
    \lambda_k>0:\qquad
    F(x_k)-F(x_{k+1})
    &\ge
    \frac{1535}{32768}
    \frac{h_k^{3/2}}{\sqrt{M/6}}
    +
    \frac12\lambda_kr_k^2.
    \label{eq:fixed-boundary-decrease}
\end{align}
\end{subequations}
Moreover,
\begin{equation}
    h_{k+1}
    \le
    \max\left\{
        \frac{2849}{4096}h_k
        +
        \lambda_kr_k,
        \epsilon
    \right\}.
    \label{eq:fixed-gradient-recurrence}
\end{equation}
\end{lemma}

\begin{proof}
Apply Lemma~\ref{lem:common-inexact-tr-step} with
\[
    a
    =
    \xi_k
    +
    2\sqrt{\frac M6}\sqrt{h_k},
    \qquad
    R
    =
    \frac{\sqrt{h_k}}{4\sqrt{M/6}},
\]
and
\[
    \xi=\xi_k,
    \qquad
    \delta_0=\delta_+
    =
    \frac{\epsilon}{8192}.
\]
Since $r_k\le R$,
\[
    \frac M6 r_k^3
    \le
    \frac14
    \sqrt{\frac M6}\sqrt{h_k}\,r_k^2.
\]
The certificate padding cancels the Hessian-error term in
\eqref{eq:common-function-bound}, giving
\[
    F(x_k)-F(x_{k+1})
    \ge
    \frac34
    \sqrt{\frac M6}\sqrt{h_k}\,r_k^2
    +
    \frac12\lambda_kr_k^2
    -
    \frac{\epsilon r_k}{8192},
\]
which proves \eqref{eq:fixed-basic-decrease}.

If $\lambda_k=0$, minimizing
\[
    \frac34\sqrt{\frac M6}\sqrt\epsilon\,r^2
    -
    \frac{\epsilon}{8192}r
\]
over $r\ge0$ gives the minimizer
\[
    r
    =
    \frac{\sqrt\epsilon}
         {12288\sqrt{M/6}}
\]
and minimum
\[
    -
    \frac{
        \epsilon^{3/2}
    }{
        201\,326\,592\sqrt{M/6}
    },
\]
which proves \eqref{eq:fixed-interior-increase}.

If $\lambda_k>0$, complementarity gives
\[
    r_k
    =
    \frac{\sqrt{h_k}}
         {4\sqrt{M/6}}.
\]
Hence
\begin{align*}
    &
    \frac34
    \sqrt{\frac M6}\sqrt{h_k}\,r_k^2
    -
    \frac{\epsilon r_k}{8192}
    \\
    &\qquad\ge
    \left(
        \frac3{64}
        -
        \frac1{32768}
    \right)
    \frac{h_k^{3/2}}{\sqrt{M/6}}
    =
    \frac{1535}{32768}
    \frac{h_k^{3/2}}{\sqrt{M/6}},
\end{align*}
which, together with the multiplier term, proves
\eqref{eq:fixed-boundary-decrease}.

Finally, the gradient estimate
\eqref{eq:common-gradient-bound} gives
\begin{align*}
    \|g_{k+1}\|
    &\le
    2\sqrt{\frac M6}\sqrt{h_k}\,r_k
    +
    2\xi_kr_k
    +
    \lambda_kr_k
    +
    \frac M2r_k^2
    +
    \frac{\epsilon}{4096}
    \\
    &\le
    \left(
        \frac12
        +
        \frac1{128}
        +
        \frac3{16}
        +
        \frac1{4096}
    \right)h_k
    +
    \lambda_kr_k
    \\
    &=
    \frac{2849}{4096}h_k+\lambda_kr_k.
\end{align*}
Here we used
\[
    2\sqrt{\frac M6}\sqrt{h_k}\,r_k
    \le
    \frac12h_k,
    \qquad
    2\xi_kr_k
    \le
    \frac1{128}h_k,
\]
\[
    \frac M2r_k^2
    \le
    \frac3{16}h_k,
    \qquad
    \frac{\epsilon}{4096}
    \le
    \frac{h_k}{4096}.
\]
Clipping below at $\epsilon$ proves
\eqref{eq:fixed-gradient-recurrence}.
\end{proof}

The negative right-hand side in
\eqref{eq:fixed-interior-increase} reflects the fact that an inexact
interior step need not decrease $F$. We next account for these small
possible increases over a finite prefix.

For any finite prefix of $N\le K$ nonterminal iterations, define
\[
    \mathcal B_N
    :=
    \{0\le k<N:\lambda_k>0\},
    \qquad
    \mathcal I_N
    :=
    \{0\le k<N:\lambda_k=0\},
\]
and, for $k\in\mathcal B_N$, set
\begin{equation}
    z_k
    :=
    \frac{\lambda_kr_k}{h_k}.
    \label{eq:fixed-z-definition}
\end{equation}

\begin{lemma}[Boundary and interior budgets]
\label{lem:fixed-budgets}
Every finite nonterminal prefix satisfies
\begin{subequations}
\label{eq:fixed-boundary-budgets}
\begin{align}
    |\mathcal B_N|
    &\le
    \frac{32768}{1535}
    \sqrt{\frac M6}
    \Delta_F\epsilon^{-3/2}
    +
    \frac1{9\,431\,040}
    |\mathcal I_N|,
    \label{eq:fixed-boundary-count}
    \\
    \sum_{k\in\mathcal B_N}z_k
    &\le
    8\sqrt{\frac M6}
    \Delta_F\epsilon^{-3/2}
    +
    \frac1{25\,165\,824}
    |\mathcal I_N|.
    \label{eq:fixed-z-budget}
\end{align}
\end{subequations}
Moreover,
\begin{equation}
    |\mathcal I_N|
    \le
    |\mathcal B_N|+1
    +
    \frac{
        \log_+(h_0/\epsilon)
        +
        \sum_{k\in\mathcal B_N}z_k
    }{
        \log(4096/2849)
    }.
    \label{eq:fixed-interior-budget}
\end{equation}
\end{lemma}

\begin{proof}
Summing \eqref{eq:fixed-step-type-decrease} and using
$F(x_N)\ge F^*$ gives
\begin{align*}
    \Delta_F
    &\ge
    \frac{
        1535\epsilon^{3/2}
    }{
        32768\sqrt{M/6}
    }
    |\mathcal B_N|
    +
    \frac12
    \sum_{k\in\mathcal B_N}\lambda_kr_k^2
    \\
    &\qquad
    -
    \frac{
        \epsilon^{3/2}
    }{
        201\,326\,592\sqrt{M/6}
    }
    |\mathcal I_N|.
\end{align*}
This directly gives
\eqref{eq:fixed-boundary-count}. For
$k\in\mathcal B_N$,
\[
    \lambda_kr_k^2
    =
    \frac{
        z_k h_k^{3/2}
    }{
        4\sqrt{M/6}
    }
    \ge
    \frac{
        z_k\epsilon^{3/2}
    }{
        4\sqrt{M/6}
    },
\]
which gives \eqref{eq:fixed-z-budget}.

For $k\in\mathcal B_N$,
\eqref{eq:fixed-gradient-recurrence} implies
\[
    \left[
        \log\frac{h_{k+1}}{h_k}
    \right]_+
    \le
    \left[
        \log\left(
            \frac{2849}{4096}+z_k
        \right)
    \right]_+
    \le
    z_k.
\]
Indeed, the left-hand side vanishes when
$z_k\le1247/4096$; above this threshold,
\[
    z-\log\left(\frac{2849}{4096}+z\right)
\]
is increasing and positive.

For an index $k\in\mathcal I_N$ with
$h_{k+1}>\epsilon$, the quantity $\log h_k$ decreases by at least
$\log(4096/2849)$.
Now suppose
$k\in\mathcal I_N$ satisfies
$h_{k+1}=\epsilon$ and $k+1<N$.
At iteration $k+1$ the gradient part of the stopping test holds, while
the procedure has not terminated. Hence
\[
    \lambda_{\min}(H_{k+1})
    <
    -\sqrt{M\epsilon}
    +
    \xi_{k+1}.
\]
Consequently,
\begin{align*}
    &
    \lambda_{\min}\!\left(
        H_{k+1}
        +
        \xi_{k+1}I
        +
        2\sqrt{\frac M6}\sqrt\epsilon\,I
    \right)
    \\
    &\qquad<
    -\sqrt{M\epsilon}
    +
    2\xi_{k+1}
    +
    2\sqrt{\frac M6}\sqrt\epsilon
    \\
    &\qquad\le
    \left(
        2+\frac1{32}-\sqrt6
    \right)
    \sqrt{\frac M6}\sqrt\epsilon
    <0.
\end{align*}
The positive-semidefiniteness condition in
Lemma~\ref{lem:tr-global-optimality} therefore forces
$\lambda_{k+1}>0$.
Thus every such floor-hitting index is followed by an index in
$\mathcal B_N$, except possibly the last prefix index. Hence
\[
    \left|
        \{k\in\mathcal I_N:h_{k+1}=\epsilon\}
    \right|
    \le
    |\mathcal B_N|+1.
\]

For the remaining indices in $\mathcal I_N$, telescoping
$\log h_k$, using $h_N\ge\epsilon$, and charging every positive
increase on $\mathcal B_N$ to $z_k$ gives
\[
    \left|
        \{k\in\mathcal I_N:h_{k+1}>\epsilon\}
    \right|
    \log(4096/2849)
    \le
    \log_+\frac{h_0}{\epsilon}
    +
    \sum_{k\in\mathcal B_N}z_k.
\]
Adding the two interior counts proves
\eqref{eq:fixed-interior-budget}.
\end{proof}


\begin{proof}[Proof of Theorem~\ref{thm:fixed-scale-utr}]
Substituting
\eqref{eq:fixed-boundary-budgets} into
\eqref{eq:fixed-interior-budget} gives
\begin{align*}
    &
    \left[
        1
        -
        \frac1{9\,431\,040}
        -
        \frac{
            1
        }{
            25\,165\,824\log(4096/2849)
        }
    \right]
    |\mathcal I_N|
    \\
    &\quad\le
    \left(
        \frac{32768}{1535}
        +
        \frac8{\log(4096/2849)}
    \right)
    \sqrt{\frac M6}
    \Delta_F\epsilon^{-3/2}
    \\
    &\qquad+
    1+
    \frac{
        \log_+(h_0/\epsilon)
    }{
        \log(4096/2849)
    }.
\end{align*}
Using
\[
    \log(4096/2849)>\frac{363}{1000}
\]
and
\[
    1
    -
    \frac1{9\,431\,040}
    -
    \frac{
        1
    }{
        25\,165\,824\log(4096/2849)
    }
    >
    \frac{9\,999\,997}{10\,000\,000},
\]
together with
\eqref{eq:fixed-boundary-count}, gives
\begin{equation}
    N
    =
    |\mathcal I_N|+|\mathcal B_N|
    \le
    65\sqrt{\frac M6}
    \Delta_F\epsilon^{-3/2}
    +
    3\log_+\frac{h_0}{\epsilon}
    +2.
    \label{eq:fixed-prefix-bound}
\end{equation}
This estimate is uniform over every finite nonterminal prefix.

If $K=\infty$, then
\eqref{eq:fixed-prefix-bound} would bound every finite
$N\le K$ by the same constant, which is impossible for arbitrarily
large $N$. Therefore $K<\infty$.
Taking $N=K$ proves
\eqref{eq:fixed-scale-complexity}.

At termination,
\eqref{eq:fixed-oracle-errors} gives
\[
    \|\nabla F(\widehat x)\|
    \le
    \frac{8193}{8192}\epsilon.
\]
By Weyl's inequality and
\eqref{eq:fixed-scale-stopping},
\[
    \lambda_{\min}(\nabla^2F(\widehat x))
    \ge
    \lambda_{\min}(H_K)-\xi_K
    \ge
    -\sqrt{M\epsilon}.
\]
This proves \eqref{eq:fixed-scale-return}.
\end{proof}


\begin{proof}[Proof of Lemma~\ref{lem:cubic-path}]
Since
\[
    r_k
    \le
    \frac{\sqrt{h_k}}
         {4\sqrt{M/6}},
\]
we have
\[
    \sqrt{h_k}
    \ge
    4\sqrt{\frac M6}\,r_k.
\]
Equation~\eqref{eq:fixed-basic-decrease} therefore implies
\[
    F(x_k)-F(x_{k+1})
    \ge
    \frac M2 r_k^3
    -
    \frac{\epsilon r_k}{8192}.
\]
Split the cubic term equally and maximize
\[
    \frac{\epsilon r}{8192}
    -
    \frac M4r^3
\]
over $r\ge0$. Its maximum is
\[
    \frac{
        \epsilon^{3/2}
    }{
        393\,216\sqrt{6M}
    }.
\]
Hence
\[
    F(x_k)-F(x_{k+1})
    \ge
    \frac M4r_k^3
    -
    \frac{
        \epsilon^{3/2}
    }{
        393\,216\sqrt{6M}
    }.
\]
Summing from $k=0$ to $K-1$ and using
$F(x_K)\ge F^*$ gives
\[
    M\sum_{k=0}^{K-1}r_k^3
    \le
    4\Delta_F
    +
    \frac{K}{98\,304\sqrt6}
    \frac{\epsilon^{3/2}}{\sqrt M},
\]
which is \eqref{eq:cubic-path}.
\end{proof}


\subsection{Oracle Calibration for the NC--SC Method}
\label{app:ncsc-oracle-calibration}

\begin{proof}[Proof of Corollary~\ref{cor:ncsc-fixed-oracle}]
Lemma~\ref{lem:inner-residual-certificate} gives
\[
    \|y-y^*(x)\|
    \le
    \frac{
        \|\nabla_yf(x,y)\|
    }{
        \mu
    }.
\]
Hence
\eqref{eq:corrected-hessian-error} gives
\[
    \|\widehat H(x,y)-\nabla^2P(x)\|
    \le
    \xi(x,y).
\]
Using
$\|\nabla_yf(x,y)\|\le\tau_{\rm out}$ and
\eqref{eq:outer-residual-target},
\[
    \|y-y^*(x)\|
    \le
    \frac{
        \sqrt{M_P\epsilon/6}
    }{
        64(1+\kappa)^2\rho
    }.
\]
Substitution into
\eqref{eq:corrected-hessian-error} gives
\[
    \xi(x,y)
    \le
    \frac1{64}
    \sqrt{\frac{M_P}{6}}\sqrt\epsilon.
\]

Similarly,
\eqref{eq:corrected-gradient-error} yields
\begin{align*}
    \|\widehat g(x,y)-\nabla P(x)\|
    &\le
    \frac{
        M_P\epsilon
    }{
        49152(1+\kappa)^3\rho
    }
    \\
    &=
    \frac{\sqrt2}{12288}
    \left(
        \frac{\kappa}{1+\kappa}
    \right)^3
    \epsilon
    <
    \frac{\epsilon}{8192}.
\end{align*}
This proves \eqref{eq:ncsc-fixed-oracle-errors}.
\end{proof}


\subsection{Inner Solver and Complexity of Algorithm~\ref{alg:ncsc-utr}}
\label{app:ncsc-complexity}


\subsubsection*{AGD realization}
\label{app:agd-realization}

\begin{proof}[Proof of Proposition~\ref{prop:agd-implementation}]
Apply the standard accelerated-gradient estimate to the
$\ell$-smooth and $\mu$-strongly convex function
\[
    \phi_x(y):=-f(x,y).
\]
The estimate
\eqref{eq:agd-contraction} follows, for example, from
\citet[Lemma~2]{wang2020improved} or
\citet[Lemma~5]{luo2022finding}.

By Lemma~\ref{lem:inner-residual-certificate},
\[
    \|y^{(0)}-y^*(x)\|
    \le
    \frac{R_{\rm in}}{\mu}.
\]
It is sufficient to obtain
\[
    \|y^{(N)}-y^*(x)\|
    \le
    \frac{\tau}{\ell},
\]
because $\ell$-smoothness then gives
\[
    \|\nabla_yf(x,y^{(N)})\|
    \le
    \tau.
\]
Combining these inequalities with
\[
    1-\frac1{\sqrt\kappa}
    \le
    \exp\!\left(-\frac1{\sqrt\kappa}\right)
\]
gives
\eqref{eq:agd-iteration-count} and
\eqref{eq:agd-module-work}.
For $\kappa=1$, one gradient step is sufficient.
\end{proof}


\subsubsection*{Complexity theorem}

\begin{proof}[Proof of Theorem~\ref{thm:ncsc-complexity}]
Let $K$ denote the number of nonterminal outer iterations and set
\[
    r_k:=\|d_k\|.
\]
By the residual contract and
Corollary~\ref{cor:ncsc-fixed-oracle},
\[
    \|g_k-\nabla P(x_k)\|
    \le
    \frac{\epsilon}{8192},
\]
and
\[
    \|H_k-\nabla^2P(x_k)\|
    \le
    \xi_k
    \le
    \frac1{64}
    \sqrt{\frac{M_P}{6}}\sqrt\epsilon.
\]
Hence the outer sequence generated by
Algorithm~\ref{alg:ncsc-utr} satisfies the fixed-scale procedure of
Section~\ref{sec:fixed-scale-utr} with
\[
    F=P,
    \qquad
    M=M_P.
\]
Theorem~\ref{thm:fixed-scale-utr} therefore gives
\eqref{eq:ncsc-return} and
\eqref{eq:ncsc-outer-explicit}.

The logarithmic term in
\eqref{eq:ncsc-outer-explicit} contains
\[
    h_0
    =
    \max\{\|g_0\|,\epsilon\}.
\]
Since
\[
    \|g_0\|
    \le
    \|\nabla P(x_0)\|
    +
    \frac{\epsilon}{8192},
\]
we have
\[
    \log_+\frac{h_0}{\epsilon}
    =
    O\!\left(
        \log_+
        \frac{\|\nabla P(x_0)\|}{\epsilon}
        +1
    \right).
\]
Together with
\[
    \sqrt{M_P}
    =
    \sqrt{4\sqrt2}\,
    \kappa^{3/2}\sqrt\rho,
\]
this proves
\eqref{eq:ncsc-outer-complexity}.

It remains to bound the total work of the inner module.
At the start of the inner call for $x_{k+1}$, joint
$\ell$-smoothness and the previous residual guarantee give
\[
    \|\nabla_yf(x_{k+1},y_k)\|
    \le
    \|\nabla_yf(x_k,y_k)\|
    +
    \ell\|x_{k+1}-x_k\|
    \le
    \tau_{\rm out}+\ell r_k.
\]
Therefore, apart from the initial call, the logarithmic term in
\eqref{eq:inner-module-work} is at most
\[
    1+
    \log_2\!\left(
        1+\frac{\ell r_k}{\tau_{\rm out}}
    \right).
\]
By \eqref{eq:outer-residual-lower},
\[
    \frac{\ell r_k}{\tau_{\rm out}}
    \le
    64\sqrt3\,
    r_k
    \sqrt{\frac{M_P}{\epsilon}}.
\]
Since
\[
    \log_2(1+t)
    =
    O(1+t^3),
    \qquad t\ge0,
\]
the logarithmic contribution of each noninitial call is
\[
    O\!\left(
        1+
        \left(
            r_k\sqrt{\frac{M_P}{\epsilon}}
        \right)^3
    \right).
\]

Applying Lemma~\ref{lem:cubic-path} with
$F=P$ and $M=M_P$ gives
\begin{align*}
    \sum_{k=0}^{K-1}
    \left(
        r_k\sqrt{\frac{M_P}{\epsilon}}
    \right)^3
    &=
    \frac{M_P^{3/2}}{\epsilon^{3/2}}
    \sum_{k=0}^{K-1}r_k^3
    \\
    &\le
    4\Delta_P\sqrt{M_P}\,
    \epsilon^{-3/2}
    +
    \frac{K}{98\,304\sqrt6}.
\end{align*}
Summing the module work bounds over all calls, adding the initial
call
\[
    O\!\left(
        C_{\mathcal S}\sqrt\kappa
        \left[
            b_{\mathcal S}
            +
            \log_+
            \frac{
                \|\nabla_yf(x_0,y_{-1})\|
            }{
                \tau_{\rm out}
            }
        \right]
    \right),
\]
and substituting the outer bound for $K$ gives
\eqref{eq:ncsc-gradient-complexity}.

Finally, if
\[
    C_{\mathcal S}=O(1),
    \qquad
    b_{\mathcal S}=O(1+\log\kappa),
\]
then
\[
    \sqrt\kappa\,\sqrt{M_P}
    =
    \Theta(\kappa^2\sqrt\rho),
\]
which gives
\eqref{eq:ncsc-gradient-leading}.
\end{proof}